% !TEX TS-program = lualatex

%%%
% src::
%     url = https://github.com/pfradin/luadraw/discussions/207#discussioncomment-15868700
%%%
\documentclass{standalone}

%%%
% Des couleurs faciles d'emploi via le package \xcolor!
%%%
\usepackage[svgnames]{xcolor}

%%%
% La bibliothèque ''luadraw'' allie une facilité d’utilisation
% à un rendu particulièrement soigné.
%%%
\usepackage[3d]{luadraw}

\begin{document}

\begin{luadraw}{name = angular-gradient-by-hand}
------
-- Par goût, nous allons utiliser la palette ''palPicnic''.
------
require 'luadraw_palettes'

------
-- Réglages globaux.
------
local Rmax = 2
local Rmin = .95
local n    = 200

------
-- Dessinons les secteurs angulaires, lesquels seront appochés
-- par des quadrilatères.
------
local graphview = graph:new{
  size = {10, 10},
  bbox = false
}

local c = -Rmin
local d = -Rmax
local a, b

local dtheta = 360 / n

local colors = getpalette(palPicnic, n)

for k = 1, n do
  a, b = c, d

  c, d = table.unpack(
    rotate({c, d}, dtheta)
  )

  graphview:Dpolyline(
    {a, b, d, c},
    true,           -- Fermons le quadrilatère.
      "color = "
    .. colors[k]
    .. ", fill = "
    .. colors[k]
  )
end

------
-- Montrons le résultat de notre oeuvre.
------
graphview:Show()
\end{luadraw}

\end{document}
